\documentclass[11pt]{article}
\usepackage[margin=1.1in]{geometry}
\usepackage{amsmath,amssymb,amsthm,mathtools,booktabs,array,tabularx,longtable}
\usepackage[dvipsnames]{xcolor}
\usepackage[colorlinks=true,linkcolor=MidnightBlue,citecolor=MidnightBlue,urlcolor=MidnightBlue]{hyperref}
\usepackage{fancyhdr}
\pagestyle{fancy}\fancyhf{}\renewcommand{\headrulewidth}{0pt}
\fancyfoot[R]{\thepage}
\fancyfoot[L]{\scriptsize\color{gray}Generated by the Transport Lab from approved parts of its store; not a deposited paper. Papers of record: \href{https://zenodo.org/communities/transport-papers}{Zenodo, Transport Papers}}
\fancypagestyle{plain}{\fancyhf{}\fancyfoot[R]{\thepage}\fancyfoot[L]{\scriptsize\color{gray}Generated by the Transport Lab from approved parts of its store; not a deposited paper. Papers of record: \href{https://zenodo.org/communities/transport-papers}{Zenodo, Transport Papers}}\renewcommand{\headrulewidth}{0pt}}
\newtheorem{theorem}{Theorem}[section]
\newtheorem{proposition}[theorem]{Proposition}
\newtheorem{lemma}[theorem]{Lemma}
\newtheorem{corollary}[theorem]{Corollary}
\newtheorem{conjecture}[theorem]{Conjecture}
\newtheorem{refutation}[theorem]{Refutation}
\newtheorem{prediction}[theorem]{Prediction}
\newtheorem{principle}[theorem]{Principle}
\newtheorem{observation}[theorem]{Observation}
\newtheorem{axiom}[theorem]{Axiom}
\theoremstyle{definition}
\newtheorem{definition}[theorem]{Definition}
\newtheorem{fact}[theorem]{Fact}
\newtheorem{assembled}[theorem]{Assembly}
\newtheorem{postulate}[theorem]{Postulate}
\newtheorem{premise}[theorem]{Premise}
\newtheorem{hypothesis}[theorem]{Hypothesis}
\newtheorem{remark}[theorem]{Remark}
\newtheorem{example}[theorem]{Example}
\theoremstyle{plain}
\setlength{\marginparwidth}{0.8in}
\newcommand{\changedmark}[1]{\marginpar{\raggedright\scriptsize\color{BrickRed}#1}}
\providecommand{\Aa}{}\renewcommand{\Aa}{\mathcal{A}}
\providecommand{\Ad}{}\renewcommand{\Ad}{\widetilde{\mathrm{Ad}}}
\providecommand{\Aev}{}\renewcommand{\Aev}{\mathcal{A}^{+}}
\providecommand{\Aod}{}\renewcommand{\Aod}{\mathcal{A}^{-}}
\providecommand{\B}{}\renewcommand{\B}{\mathrm{B}}
\providecommand{\C}{}\renewcommand{\C}{\mathbb{C}}
\providecommand{\CAR}{}\renewcommand{\CAR}{\mathrm{CAR}}
\providecommand{\CS}{}\renewcommand{\CS}{{\mathbb{S}^{+}}}
\providecommand{\CSm}{}\renewcommand{\CSm}{{\mathbb{S}^{-}}}
\providecommand{\Cl}{}\renewcommand{\Cl}{\mathrm{Cl}}
\providecommand{\Con}{}\renewcommand{\Con}{\mathcal{C}}
\providecommand{\Dl}{}\renewcommand{\Dl}{\Delta}
\providecommand{\EN}{}\renewcommand{\EN}{E_{N}}
\providecommand{\End}{}\renewcommand{\End}{\mathrm{End}}
\providecommand{\F}{}\renewcommand{\F}{\mathbb{F}}
\providecommand{\G}{}\renewcommand{\G}{\Gamma}
\providecommand{\HH}{}\renewcommand{\HH}{\mathbb{H}}
\providecommand{\IS}{}\renewcommand{\IS}{I\sigma_3}
\providecommand{\J}{}\renewcommand{\J}{\mathbf{J}}
\providecommand{\LH}{}\renewcommand{\LH}{L_{H}}
\providecommand{\Mat}{}\renewcommand{\Mat}{\mathrm{M}}
\providecommand{\Mo}{}\renewcommand{\Mo}{M^{\circ}}
\providecommand{\Mp}{}\renewcommand{\Mp}{M^{+}}
\providecommand{\Msun}{}\renewcommand{\Msun}{M_{\odot}}
\providecommand{\Mt}{}\renewcommand{\Mt}{\widetilde{M}}
\providecommand{\Ncone}{}\renewcommand{\Ncone}{\mathcal{N}}
\providecommand{\Neff}{}\renewcommand{\Neff}{N_{\mathrm{eff}}}
\providecommand{\PCS}{}\renewcommand{\PCS}{\mathbb{P}{\mathbb{S}^{+}}}
\providecommand{\PT}{}\renewcommand{\PT}{\mathbb{PT}}
\providecommand{\Pin}{}\renewcommand{\Pin}{\mathrm{Pin}}
\providecommand{\Q}{}\renewcommand{\Q}{\mathbb{Q}}
\providecommand{\R}{}\renewcommand{\R}{\mathbb{R}}
\providecommand{\Ric}{}\renewcommand{\Ric}{\mathrm{Ric}}
\providecommand{\SBH}{}\renewcommand{\SBH}{S_{\mathrm{BH}}}
\providecommand{\Sc}{}\renewcommand{\Sc}{\mathbb{S}}
\providecommand{\Sch}{}\renewcommand{\Sch}{\mathrm{P}}
\providecommand{\Sig}{}\renewcommand{\Sig}{\Sigma}
\providecommand{\Som}{}\renewcommand{\Som}{S_\omega}
\providecommand{\Sp}{}\renewcommand{\Sp}{\mathbb{S}}
\providecommand{\Sph}{}\renewcommand{\Sph}{\mathrm{S}}
\providecommand{\Spi}{}\renewcommand{\Spi}{S_\pi}
\providecommand{\Spin}{}\renewcommand{\Spin}{\mathrm{Spin}}
\providecommand{\Sth}{}\renewcommand{\Sth}{S_{\mathrm{th}}}
\providecommand{\T}{}\renewcommand{\T}{\mathbb{T}}
\providecommand{\TH}{}\renewcommand{\TH}{T_{\mathrm{H}}}
\providecommand{\Tr}{}\renewcommand{\Tr}{\operatorname{Tr}}
\providecommand{\UHF}{}\renewcommand{\UHF}{\mathrm{UHF}}
\providecommand{\Z}{}\renewcommand{\Z}{\mathbb{Z}}
\providecommand{\Zc}{}\renewcommand{\Zc}{\mathrm{Z}}
\providecommand{\cmg}{}\renewcommand{\cmg}{\mathrm{cm^{2}\,g^{-1}}}
\providecommand{\dS}{}\renewcommand{\dS}{\mathrm{dS}}
\providecommand{\dsc}{}\renewcommand{\dsc}{\mathrm{dsc}}
\providecommand{\dtPG}{}\renewcommand{\dtPG}{dt}
\providecommand{\dtS}{}\renewcommand{\dtS}{dt_{\mathrm{S}}}
\providecommand{\ellv}{}\renewcommand{\ellv}{\ell}
\providecommand{\er}{}\renewcommand{\er}{e_r}
\providecommand{\et}{}\renewcommand{\et}{e_t}
\providecommand{\ev}{}\renewcommand{\ev}{\mathrm{ev}}
\providecommand{\gl}{}\renewcommand{\gl}{\mathfrak{gl}}
\providecommand{\hatV}{}\renewcommand{\hatV}{\widehat{V}}
\providecommand{\hb}{}\renewcommand{\hb}{\bar{h}}
\providecommand{\hf}{}\renewcommand{\hf}{\hat{f}}
\providecommand{\id}{}\renewcommand{\id}{\mathrm{id}}
\providecommand{\kpc}{}\renewcommand{\kpc}{\mathrm{kpc}}
\providecommand{\lP}{}\renewcommand{\lP}{\ell_{P}}
\providecommand{\mD}{}\renewcommand{\mD}{\mathbf{m}_{D}}
\providecommand{\mb}{}\renewcommand{\mb}{\bar m}
\providecommand{\mnu}{}\renewcommand{\mnu}{\Sigma m_{\nu}}
\providecommand{\nInf}{}\renewcommand{\nInf}{n_\infty}
\providecommand{\nO}{}\renewcommand{\nO}{n_0}
\providecommand{\ninf}{}\renewcommand{\ninf}{n_{\infty}}
\providecommand{\nz}{}\renewcommand{\nz}{n_{0}}
\providecommand{\proj}{}\renewcommand{\proj}{\operatorname{proj}}
\providecommand{\pv}{}\renewcommand{\pv}{\pi}
\providecommand{\rad}{}\renewcommand{\rad}{\mathrm{rad}}
\providecommand{\sgn}{}\renewcommand{\sgn}{\operatorname{sgn}}
\providecommand{\shA}{}\renewcommand{\shA}{\mathcal{A}}
\providecommand{\shC}{}\renewcommand{\shC}{\mathcal{C}}
\providecommand{\shE}{}\renewcommand{\shE}{\mathcal{E}}
\providecommand{\shJ}{}\renewcommand{\shJ}{\mathcal{J}}
\providecommand{\shK}{}\renewcommand{\shK}{\mathcal{K}}
\providecommand{\shL}{}\renewcommand{\shL}{L}
\providecommand{\shW}{}\renewcommand{\shW}{\mathcal{W}}
\providecommand{\sll}{}\renewcommand{\sll}{\mathfrak{sl}}
\providecommand{\so}{}\renewcommand{\so}{\mathfrak{so}}
\providecommand{\su}{}\renewcommand{\su}{\mathfrak{su}}
\providecommand{\supp}{}\renewcommand{\supp}{\operatorname{supp}}
\providecommand{\tr}{}\renewcommand{\tr}{\operatorname{tr}}
\providecommand{\uu}{}\renewcommand{\uu}{\mathfrak{u}}
\providecommand{\br}{}\renewcommand{\br}[1]{\langle #1\rangle}
\providecommand{\bra}{}\renewcommand{\bra}[1]{\langle #1|}
\providecommand{\file}{}\renewcommand{\file}[1]{\texttt{#1}}
\providecommand{\hp}{}\renewcommand{\hp}[2]{( #1 \,|\, #2 )}
\providecommand{\ip}{}\renewcommand{\ip}[2]{\langle #1, #2\rangle}
\providecommand{\ket}{}\renewcommand{\ket}[1]{|#1\rangle}
\providecommand{\rel}{}\renewcommand{\rel}[2]{D\bigl(#1 \,\|\, #2\bigr)}
\providecommand{\rev}{}\renewcommand{\rev}[1]{\widetilde{#1}}
\providecommand{\slashed}{}\renewcommand{\slashed}[1]{{#1}\mkern-10mu/}
\title{\bfseries Rule 2 Exercised Physically: Kaluza-Klein directions, planar fermions, and the commuting square: Companion note to T17: closing Direction 7 of the forward ledger}
\author{Leon Butler\\ \normalsize\href{mailto:lbutler.math@gmail.com}{\texttt{lbutler.math@gmail.com}}\\ \normalsize ORCID: \href{https://orcid.org/0009-0006-4449-5617}{0009-0006-4449-5617}}
\date{Generated by the Transport Lab, 27 October 2026}
\begin{document}
\maketitle
\begin{abstract}
\normalsize
The capstone's forward ledger leaves one direction unowned: a physical system whose natural algebra carries a central splitting, so that Rule 2 of the transport principle, ideal fusion by one adjoined generator, is exercised by nature rather than by the classification table. This note supplies two, and the theorem behind both. Of the two one-dimensional extensions of an even-dimensional nondegenerate Clifford algebra exactly one carries the type-C obstruction, spacelike when $p - q \equiv 2 \pmod 4$ and timelike when $p - q \equiv 0$, the other having complex type; each ideal of the obstructed extension is a copy of the base; and the Rule 2 repair by a generator of the opposite sign is exactly the Rule 1 ascent $\Cl(p{+}1,q{+}1)$: the two adjoined directions span its hyperbolic plane and carry its ladder pair, and the even half that Rule 2 amplifies is the base algebra itself. Applied to the physical lineage: the five-dimensional Kaluza-Klein algebra $\Cl(1,4)$ splits, its two ideals are the two inequivalent five-dimensional Dirac representations, the fifth-direction reflection is the outer exchange, and the fusion is the conformal algebra $\Cl(2,4)$ or its spatial twin $\Cl(1,5)$, the bit the algebra forgets being whether the parity becomes a boost or a rotation. Applied to the plane: $\Cl(2,1)$ splits into the two two-component representations, the four-component fermion of planar electrodynamics and of graphene is the fused algebra, and the parity anomaly is the obstruction's physical face. Twenty-four exact checks verify every identity.
\end{abstract}
\medskip
\noindent\textbf{MSC 2020.} Primary 15A66; Secondary 81R25, 83E15, 81T25.\\
\textbf{Keywords.} Clifford algebras; central idempotents; Kaluza-Klein; parity anomaly; graphene; conformal algebra; transport principle.
\newpage
\noindent \textbf{Authorship and disclosure. }The founding idea of the series is the author's: that division by zero in a Clifford algebra should be read not as a prohibition but as a transport into a higher-dimensional algebra, and that the transported division need not be accompanied by an inverse operation. The direction of investigation and the critical review of the results are likewise the author's. The mathematical development and the drafting of the text were carried out in extended collaboration with Claude Fable 5 (Anthropic), an AI system.

\noindent\fbox{\parbox{\dimexpr\linewidth-2\fboxsep-2\fboxrule}{\textbf{This is not a deposited paper.} It was generated by the Transport Lab on 27 October 2026 and assembled from approved parts of its store. It follows results across several papers of the series; each is stated with its proof and credited to the paper that first proved it. The official versions of those papers are on Zenodo, in the Transport Papers community, \url{https://zenodo.org/communities/transport-papers}.}}
\newpage

\section{Setting, and what is imported}

\phantomsection\label{c:note.rule2.setting-setting-what-imported}
The transport principle sorts the zero divisors of a real Clifford algebra into three types and repairs each by a minimal ascent \cite{ButlerT1}. Type C is the deepest: not a direction without a partner but a splitting of the algebra itself into two ideals, along the central idempotents $\varepsilon_\pm = \tfrac12(1 \pm \omega)$ built from the volume element, which no inner automorphism can exchange. Its repair, Rule 2, adjoins one generator, fuses the ideals into a simple algebra, and renders inner the very symmetry whose absence defined the obstruction. The capstone \cite{ButlerT17} records that every other direction of the programme has an owner among the series' papers, and that this one does not: no paper of the series exhibits a physical system whose natural algebra is split, with Kaluza-Klein extensions the named candidates. This note takes the row.

The claim is modest and the standard is the series'. What is proved here is algebra: where the obstruction sits among the one-dimensional extensions of an even-dimensional algebra, what its ideals are, and what its repair is. What is physical is imported as Facts with provenance: that odd-dimensional gamma matrices come in two inequivalent representations, that a single planar fermion breaks parity through an induced Chern-Simons term while a doubled one does not, that graphene's low-energy theory is that doubled fermion, that a five-dimensional fermion shifts a Chern-Simons level by half the sign of its mass. The note's contribution is the identification: those two representations are the two ideals, the parity that relates them is the outer exchange, the doubling is the fusion, and the fused algebra of the physical lineage is the conformal algebra the series has used since its root paper. the section ``What has been exercised, and what has not'' says exactly what has been exercised and what has not.

Notation follows the root paper: $\Cl(p,q)$ has $p$ generators of square $+1$ and $q$ of square $-1$, $n = p + q$, $d = p - q \bmod 8$ is the clock residue, $\omega = e_1 \cdots e_n$ is the volume element, $\Cl^0$ is the even subalgebra, and $\hat{a}$ is the grade involution. The series labels this world's algebra $\Cl(1,3)$ (time of square $+1$), residue $6$; the dark-matter paper explains why the labelling is a convention and what changes hands under its reversal \cite{ButlerT14}, and Corollary \ref{c:rule2.odd-lorentzian-dimension-exactly-one-real} below is where that reversal bites. Below, ``spacelike'' and ``timelike'' name generators of square $-1$ and $+1$ respectively, the series' convention; in the mirror's convention the names swap and the mathematics does not.

\begin{fact}[Imported: classification, the centre, and type C (external ingredients)]\label{c:rule2.imported-classification-centre-type-c-external}
The results from outside the series on which \ref{c:rule2.imported-classification-centre-type-c} draws, as cited there: M. F. Atiyah, R. Bott, A. Shapiro, Clifford modules, Topology 3, Suppl. 1 (1964) 3--38.; P. Lounesto, Clifford Algebras and Spinors, 2nd ed., London Math. Soc. Lecture Note Ser. 286, Cambridge University Press (2001)..
\emph{Source:} M. F. Atiyah, R. Bott, A. Shapiro, Clifford modules, Topology 3, Suppl. 1 (1964) 3--38.; P. Lounesto, Clifford Algebras and Spinors, 2nd ed., London Math. Soc. Lecture Note Ser. 286, Cambridge University Press (2001)..
\end{fact}

\begin{assembled}[Imported: classification, the centre, and type C]\label{c:rule2.imported-classification-centre-type-c}
$\Cl(p,q)$ is simple for $n$ even, with centre $\R$. For $n$ odd the volume element is central, $\omega^2 = (-1)^{n(n-1)/2}(-1)^q$, and the centre is $\R[\omega]$: isomorphic to $\R \oplus \R$ when $\omega^2 = +1$ and to $\C$ when $\omega^2 = -1$. $\Cl(p,q)$ contains a nontrivial central idempotent if and only if $p - q \equiv 1 \pmod 4$, and then the nontrivial central idempotents are exactly $\varepsilon_\pm = \tfrac12(1 \pm \omega)$, exchanged by the grade involution, which is outer on the split algebra: no inner automorphism exchanges the two ideals $\varepsilon_\pm \Cl(p,q)$, each of which is simple. \emph{Provenance:} \cite[Thm.~2.5, Lemmas 2.6 and 2.7, Thm.~3.6]{ButlerT1}; the table itself is classical \cite{ABS1964, Lounesto2001}.
\end{assembled}
\begin{proof}
$\Cl(p,q)$ is simple for $n$ even, with centre $\R$. For $n$ odd the volume element is central, $\omega^2 = (-1)^{n(n-1)/2}(-1)^q$, and the centre is $\R[\omega]$: isomorphic to $\R \oplus \R$ when $\omega^2 = +1$ and to $\C$ when $\omega^2 = -1$. $\Cl(p,q)$ contains a nontrivial central idempotent if and only if $p - q \equiv 1 \pmod 4$, and then the nontrivial central idempotents are exactly $\varepsilon_\pm = \tfrac12(1 \pm \omega)$, exchanged by the grade involution, which is outer on the split algebra: no inner automorphism exchanges the two ideals $\varepsilon_\pm \Cl(p,q)$, each of which is simple. \emph{Provenance:} \cite[Thm.~2.5, Lemmas 2.6 and 2.7, Thm.~3.6]{ButlerT1} (\ref{c:transport.centre}, \ref{c:transport.existence-classification-type-c}, \ref{c:transport.classification-real-clifford-algebras}, \ref{c:transport.sign}); the table itself is classical \cite{ABS1964, Lounesto2001}.

This restates \ref{c:rule2.imported-classification-centre-type-c-external}.

\end{proof}

\begin{assembled}[Imported: Rule 2, the matrix units, and the forgotten bit]\label{c:rule2.imported-rule-2-matrix-units-forgotten}
Let $p - q \equiv 1 \pmod 4$, $A = \Cl(p,q)$, and let $\nu$ be one adjoined generator orthogonal to the generating space with $\nu^2 = \gamma \in \{\pm 1\}$, so that $B_\gamma = \Cl(p{+}1,q)$ or $\Cl(p,q{+}1)$. Then: conjugation by $\nu$ restricts on $A$ to the grade involution, $\nu a \nu^{-1} = \hat{a}$; the map $a \mapsto \varepsilon_+ a$ is an isomorphism $\Cl^0(p,q) \to \varepsilon_+ A$; the elements $E_{11} = \varepsilon_+$, $E_{22} = \varepsilon_-$, $E_{12} = \varepsilon_+\nu$, $E_{21} = \gamma\varepsilon_-\nu$ form a full $2 \times 2$ system of matrix units in $B_\gamma$; hence $B_\gamma \cong \Mat_2(\Cl^0(p,q))$ is simple, the images of $\varepsilon_\pm$ are a conjugate pair of non-central idempotents exchanged by $\Ad_\nu$, and the two ascents are isomorphic as algebras though non-isometric as quadratic data: the sign $\gamma$ is a geometric datum with no algebraic shadow. One nondegenerate generator is minimal, the null generator never fuses, and fusion fires at most once in any transport history. \emph{Provenance:} \cite[Prop.~5.2, Lemma 5.3, Lemma 5.4, Thm.~5.5, Thm.~5.7, Remark 5.8, Prop.~7.4]{ButlerT1}.
\end{assembled}
\begin{proof}
Let $p - q \equiv 1 \pmod 4$, $A = \Cl(p,q)$, and let $\nu$ be one adjoined generator orthogonal to the generating space with $\nu^2 = \gamma \in \{\pm 1\}$, so that $B_\gamma = \Cl(p{+}1,q)$ or $\Cl(p,q{+}1)$. Then: conjugation by $\nu$ restricts on $A$ to the grade involution, $\nu a \nu^{-1} = \hat{a}$; the map $a \mapsto \varepsilon_+ a$ is an isomorphism $\Cl^0(p,q) \to \varepsilon_+ A$; the elements $E_{11} = \varepsilon_+$, $E_{22} = \varepsilon_-$, $E_{12} = \varepsilon_+\nu$, $E_{21} = \gamma\varepsilon_-\nu$ form a full $2 \times 2$ system of matrix units in $B_\gamma$; hence $B_\gamma \cong \Mat_2(\Cl^0(p,q))$ is simple, the images of $\varepsilon_\pm$ are a conjugate pair of non-central idempotents exchanged by $\Ad_\nu$, and the two ascents are isomorphic as algebras though non-isometric as quadratic data: the sign $\gamma$ is a geometric datum with no algebraic shadow. One nondegenerate generator is minimal, the null generator never fuses, and fusion fires at most once in any transport history. \emph{Provenance:} \cite[Prop.~5.2, Lemma 5.3, Lemma 5.4, Thm.~5.5, Thm.~5.7, Remark 5.8, Prop.~7.4]{ButlerT1} (\ref{c:transport.fusion-occurs-most-once}, \ref{c:transport.rule-2-ideal-fusion-sharp-form}, \ref{c:transport.corner-even-subalgebra}, \ref{c:transport.matrix-units-peirce-isomorphism}, \ref{c:transport.one-generator-minimal-minimal-completions-classified}, \ref{c:transport.ascent-renders-grade-involution-inner}).

\end{proof}

\begin{assembled}[Imported: the completion of the spacetime algebra, and its spinors]\label{c:rule2.imported-completion-spacetime-algebra-spinors}
The forced completion of the spacetime algebra is $\Cl(1,3) \hookrightarrow \Cl(2,4)$, and the complexified $\Cl(2,4)$ acts irreducibly on the Fock space of three fermionic ladder pairs, whose chirality halves are the conformal spinor space $\CS$ (Penrose's twistor space; in the real algebra it is the minimal left ideal $\Cl^+(2,4)P$, whose complex unit is $I_6$) and its conjugate dual. \emph{Provenance:} \cite[Thms.~4.2 and 4.13]{ButlerT1}; \cite[Thms.~3.1 and 3.4]{ButlerT9}.
\end{assembled}
\begin{proof}
The forced completion of the spacetime algebra is $\Cl(1,3) \hookrightarrow \Cl(2,4)$, and the complexified $\Cl(2,4)$ acts irreducibly on the Fock space of three fermionic ladder pairs, whose chirality halves are the conformal spinor space $\CS$ (Penrose's twistor space; in the real algebra it is the minimal left ideal $\Cl^+(2,4)P$, whose complex unit is $I_6$) and its conjugate dual. \emph{Provenance:} \cite[Thms.~4.2 and 4.13]{ButlerT1} (\ref{c:transport.no-go-level-elements}, \ref{c:transport.minimality}); \cite[Thms.~3.1 and 3.4]{ButlerT9} (\ref{c:correspondence.fock-model-complexified-ascended-algebra}, \ref{c:correspondence.spinor-identification}).

\end{proof}

\section{One-dimensional extensions of an even-dimensional algebra}

\phantomsection\label{c:note.rule2.setting-one-dimensional-extensions-even-dimensional}
The physical algebras of the series are even-dimensional and simple. A Kaluza-Klein extension adjoins one direction, and the algebra becomes odd-dimensional, where the volume element is central and the type-C question is live. The first result places the obstruction exactly.

\begin{lemma}[The even subalgebra of a one-dimensional extension]\label{c:rule2.even-subalgebra-one-dimensional-extension}
Let $\nu$ be a generator adjoined orthogonally to $\Cl(p,q)$. If $\nu^2 = -1$ then $e_i \mapsto e_i\nu$ extends to an isomorphism $\Cl(p,q) \cong \Cl^0(p,q{+}1)$; if $\nu^2 = +1$ then the same map gives $\Cl(q,p) \cong \Cl^0(p{+}1,q)$.
\end{lemma}
\begin{proof}
The images are even. Their squares are $(e_i\nu)^2 = e_i\nu e_i\nu = -e_i^2\nu^2$, which is $e_i^2$ for $\nu^2 = -1$ and $-e_i^2$ for $\nu^2 = +1$; for $i \neq j$, $(e_i\nu)(e_j\nu) = -e_ie_j\nu^2$ and $(e_j\nu)(e_i\nu) = -e_je_i\nu^2 = e_ie_j\nu^2$, so distinct images anticommute. By the universal property (\ref{c:std-algebra.existence-uniqueness}) the assignment extends to a unital homomorphism from $\Cl(p,q)$ (respectively $\Cl(q,p)$) into the even subalgebra. It is surjective: for $i \neq j$, $(e_i\nu)(e_j\nu) = -\nu^2 e_ie_j = \mp e_ie_j$, so every even blade of the extension not involving $\nu$ is, up to sign, a product of images, and every even blade involving $\nu$ is $e_I\nu$ with $|I|$ odd, which is $\pm(e_{i_1}\nu)(e_{i_2}\nu)(e_{i_3}\nu)\cdots$ after pairing all but the last factor. Both sides have dimension $2^n$, so the surjection is an isomorphism. The script confirms the relations, the evenness, and the dimension on $(1,3)$ and $(2,0)$ with both signs of $\nu^2$.

\end{proof}

\begin{theorem}[Where the obstruction sits, and the commuting square]\label{c:rule2.where-obstruction-sits-commuting-square}
Let $A = \Cl(p,q)$ with $n = p + q$ even, and let $A_s = \Cl(p,q{+}1)$ and $A_t = \Cl(p{+}1,q)$ be its spacelike and timelike one-dimensional extensions.
\begin{enumerate}
\item[(a)] Exactly one of $A_s$, $A_t$ carries the type-C obstruction: $A_s$ if $p - q \equiv 2 \pmod 4$, $A_t$ if $p - q \equiv 0 \pmod 4$. The other has complex type: its volume element is central with square $-1$, its centre is $\C$, and it is simple.
\item[(b)] Each of the two ideals of the obstructed extension is a copy of $A$: $\varepsilon_\pm A_s \cong \Cl^0(p,q{+}1) \cong A$ in the first case, $\varepsilon_\pm A_t \cong \Cl^0(p{+}1,q) \cong \Cl(q,p) \cong A$ in the second (the last isomorphism because the residues $0$ and $4$ are fixed by negation mod $8$).
\item[(c)] The Rule 2 repair of the obstructed extension by a generator of the opposite sign is $\Cl(p{+}1,q{+}1)$, the Rule 1 ascent of $A$: the two adjoined generators $e_{n+1}$ (the extension) and $\nu$ (the repair) span a hyperbolic plane, the null vectors $f = \tfrac12(\nu \pm e_{n+1})$, $f^* = \tfrac12(\nu \mp e_{n+1})$ (signs chosen so that each is null) are the ladder pair of Fact \ref{c:transport.rule-1-morita-trivial}, and the two amplifications agree: $\Mat_2(\Cl^0(\text{extension})) = \Mat_2(A)$. The residue returns to $d$.
\item[(d)] The Rule 2 repair by a generator of the same sign is $\Cl(p,q{+}2)$ or $\Cl(p{+}2,q)$, isomorphic to $\Cl(p{+}1,q{+}1)$ as an algebra, at residue $d \mp 2$; the adjoined plane is definite, carries no null vector, and admits no Rule 1 reading.
\end{enumerate}
\end{theorem}
\begin{proof}
(a) By Fact \ref{c:rule2.imported-classification-centre-type-c}, an odd-dimensional algebra is type C exactly when $p' - q' \equiv 1 \pmod 4$ and complex type exactly when $p' - q' \equiv 3 \pmod 4$; these exhaust the odd residues. For $A_s$, $p' - q' = p - q - 1$, which is $\equiv 1$ iff $p - q \equiv 2$ and $\equiv 3$ iff $p - q \equiv 0 \pmod 4$. For $A_t$, $p' - q' = p - q + 1$, which is $\equiv 1$ iff $p - q \equiv 0$ and $\equiv 3$ iff $p - q \equiv 2$. Since $n$ is even, $p - q$ is even, so exactly one case obtains for each extension and the two extensions fall on opposite sides. Simplicity of the complex-type algebra: its centre $\C$ is a field, and a semisimple algebra with a field for centre has one Wedderburn block.

(b) The corner isomorphism $\varepsilon_+ A' \cong \Cl^0(A')$ is Fact \ref{c:rule2.imported-rule-2-matrix-units-forgotten}, and $\varepsilon_- A' \cong \varepsilon_+ A'$ through the grade involution, which exchanges the ideals. Lemma \ref{c:rule2.even-subalgebra-one-dimensional-extension} identifies $\Cl^0(p,q{+}1) \cong \Cl(p,q)$ and $\Cl^0(p{+}1,q) \cong \Cl(q,p)$. In the timelike case, $p - q \equiv 0 \pmod 4$ gives $q - p \equiv -(p-q) \pmod 8$, and the classification of Fact \ref{c:rule2.imported-classification-centre-type-c} assigns $\Cl(q,p)$ and $\Cl(p,q)$ the same type and dimension when $p - q \equiv 0 \pmod 4$ (residues $0$ and $4$ are each fixed by negation mod 8), so $\Cl(q,p) \cong \Cl(p,q)$ there (the script exercises the timelike form of Lemma \ref{c:rule2.even-subalgebra-one-dimensional-extension} on $(1,3)$ and $(2,0)$).

(c) The spacelike case: the repair generator $\nu$ has $\nu^2 = +1$ and the extension generator $e_{n+1}$ has square $-1$, so the Gram matrix of the pair is $\mathrm{diag}(+1,-1)$, of signature $(1,1)$: a hyperbolic plane. Then $f = \tfrac12(\nu + e_{n+1})$ satisfies $f^2 = \tfrac14(\nu^2 + e_{n+1}^2 + \nu e_{n+1} + e_{n+1}\nu) = 0$, likewise $f^* = \tfrac12(\nu - e_{n+1})$, and $ff^* + f^*f = \tfrac14\bigl((\nu + e_{n+1})(\nu - e_{n+1}) + (\nu - e_{n+1})(\nu + e_{n+1})\bigr) = \tfrac12(\nu^2 - e_{n+1}^2) = 1$. The quadratic space is $\R^{p,q} \perp \langle e_{n+1}\rangle \perp \langle \nu \rangle = \R^{p+1,q+1}$ whichever order the two directions are adjoined in, so the algebra is $\Cl(p{+}1,q{+}1)$ on the nose, and Fact \ref{c:transport.rule-1-morita-trivial} reads it as $\Mat_2(A)$ through the ladder pair, while Fact \ref{c:rule2.imported-rule-2-matrix-units-forgotten} reads it as $\Mat_2(\Cl^0(A_s))$ through $\varepsilon_\pm$ and $\nu$; by (b) the two coefficient algebras are the same algebra. The residue of $\Cl(p{+}1,q{+}1)$ is $p - q$. The timelike case is the same computation with the signs of $\nu^2$ and $e_{n+1}^2$ exchanged and $f = \tfrac12(e_{n+1} + \nu)$, $f^* = \tfrac12(e_{n+1} - \nu)$.

(d) Fact \ref{c:rule2.imported-rule-2-matrix-units-forgotten} gives the isomorphism of the two repairs and the residues $d' \pm 1$ of the fused algebras with $d' = d \mp 1$ the residue of the extension; the same-sign repair has residue $d \mp 2$. The adjoined plane has Gram matrix $\pm\,\mathrm{diag}(1,1)$, definite, so $(a e_{n+1} + b\nu)^2 = \pm(a^2 + b^2) \neq 0$ for $(a,b) \neq 0$: no null vector, hence no ladder pair, hence no reading of the repair as a $(1,1)$ ascent of anything.

The script verifies (a) on every even-dimensional $(p,q)$ with $2 \leq n \leq 6$ by computing the volume element's centrality and square directly, and verifies (c) and (d) on the same grid: the obstructed extension is type C, the opposite-sign repair sits at residue $d$ with a hyperbolic adjoined plane carrying a ladder pair, and the same-sign repair sits two residues away with a definite plane.

\end{proof}

\begin{corollary}[Which residues obstruct which way]\label{c:rule2.which-residues-obstruct-which-way}
The even residues whose spacelike extension is type C are $\{2, 6\}$; those whose timelike extension is type C are $\{0, 4\}$. In particular every even residue has exactly one obstructing one-dimensional extension. The physical lineage sits at residue $6$ and its Rule 1 ascents never leave it (Fact \ref{c:transport.rule-1-morita-trivial}); by Theorem \ref{c:rule2.where-obstruction-sits-commuting-square} it leaves residue $6$ only through a spacelike Kaluza-Klein direction, into the obstruction at residue $5$, and the opposite-sign repair returns it to $6$ as $\Mat_2$ of where it started. The same holds at the mirror residue $2$, with spacelike read as square $-1$ in both towers: in T14's causal convention, where the mirror's time is its minus sign, that extension is the mirror's second time, and the mirror's own Kaluza-Klein direction gives the unobstructed $\Cl(4,1)$ (Proposition \emph{The mirror's extensions under the causal convention}).
\end{corollary}
\begin{proof}
Residues $2$ and $6$ are the even residues $\equiv 2 \pmod 4$; residues $0$ and $4$ are those $\equiv 0$. The rest is Theorem \ref{c:rule2.where-obstruction-sits-commuting-square}(a) and (c) read at $d = 6$ and $d = 2$, with Rule 1 ascents fixing the residue by \ref{c:transport.rule-1-morita-trivial}. The script lists both sets from the grid.

\end{proof}

\phantomsection\label{c:note.rule2.setting-one-dimensional-extensions-even-dimensional-2}
The dark-matter paper's oriented and causal residues are exactly $\{2, 6\}$ \cite[Cor.~3.11]{ButlerT14}. That the towers carrying matter are exactly the towers whose spacelike extensions obstruct is, at this level, arithmetic ($2 \bmod 4$ twice over); the note records the coincidence without building on it.

\begin{corollary}[Odd Lorentzian dimension: exactly one real form is split]\label{c:rule2.odd-lorentzian-dimension-exactly-one-real}
Let $n$ be odd and consider the two real forms of the Lorentzian algebra in dimension $n$, $\Cl(1,n{-}1)$ (time of square $+1$) and $\Cl(n{-}1,1)$ (time of square $-1$). Then $\Cl(1,n{-}1)$ is type C if and only if $n \equiv 1 \pmod 4$, $\Cl(n{-}1,1)$ is type C if and only if $n \equiv 3 \pmod 4$, and in each case the other real form has complex type. The two real forms have the same complexification, which is split for every odd $n$: over $\C$ the element $\omega' = \omega$ or $i\omega$, whichever has square $+1$, is central, and $\tfrac12(1 \pm \omega')$ are its central idempotents. Thus the obstruction is present in every odd Lorentzian dimension; the real form decides whether it is visible as a splitting or hidden as a central complex structure.
\end{corollary}
\begin{proof}
$\Cl(1,n{-}1)$ is the spacelike extension of $\Cl(1,n{-}2)$, an even-dimensional algebra with $p - q = 3 - n$, which is $\equiv 2 \pmod 4$ iff $n \equiv 1 \pmod 4$ and $\equiv 0$ iff $n \equiv 3$; Theorem \ref{c:rule2.where-obstruction-sits-commuting-square}(a) gives the first claim and the complex type of the other case. $\Cl(n{-}1,1)$ is the extension of $\Cl(n{-}2,1)$, $p - q = n - 3$, by a generator of square $+1$, obstructed iff $n - 3 \equiv 0$, i.e.\ $n \equiv 3 \pmod 4$, complex type otherwise. For the complexification: over $\C$ the substitution $e_i \mapsto ie_i$ on the generators of square $-1$ turns the relations of $\Cl(p,q)$ into those of $\Cl(n,0)$, so $\Cl(p,q) \otimes \C$ depends only on $n$ (\ref{c:std-algebra.complexification}); the volume element remains central, and if $\omega^2 = -1$ then $(i\omega)^2 = +1$, so $\omega'$ as defined is a central involution and $\tfrac12(1 \pm \omega')$ are orthogonal central idempotents summing to $1$. The script checks the alternation, with its complex-type complement, in dimensions $3, 5, 7, 9, 11$.

\end{proof}

\begin{remark}[The three dimensions that matter below]\label{c:note.rule2.three-dimensions-matter-below}
$n = 5$: $\Cl(1,4)$ is split, $\Cl(4,1)$ complex type. $n = 3$: $\Cl(2,1)$ is split, $\Cl(1,2)$ complex type. $n = 11$: $\Cl(10,1)$ is split, $\Cl(1,10)$ complex type. In the series' own convention the five-dimensional Kaluza-Klein extension of this world is the split one; the planar and eleven-dimensional cases are split in the mirror's convention. Physics done with complex gamma matrices sees the splitting in every odd dimension, as the corollary says it must; what the real form governs is the reality of the spinors, Remark \ref{c:note.rule2.eleven-dimensions-reality-spinors}.
\end{remark}

\section{The five-dimensional Kaluza-Klein algebra}

\phantomsection\label{c:note.rule2.setting-five-dimensional-kaluza-klein-algebra}
Take the physical algebra $\Cl(1,3)$ with generators $e_0$ (square $+1$) and $e_1, e_2, e_3$ (square $-1$), spacetime pseudoscalar $\gamma = e_0e_1e_2e_3$ with $\gamma^2 = -1$, and adjoin one spacelike direction $e_4$, $e_4^2 = -1$: the algebra $\Cl(1,4)$ of a five-dimensional spacetime with one extra spatial dimension, which is what a Kaluza-Klein reduction on a circle sees pointwise.

\begin{fact}[Imported: gamma matrices in odd dimension]\label{c:rule2.imported-gamma-matrices-odd-dimension}
In odd spacetime dimension $n$ the complex Clifford algebra has exactly two inequivalent irreducible representations, each of dimension $2^{(n-1)/2}$, related by $\Gamma^M \mapsto -\Gamma^M$ for all $M$ and distinguished by the scalar value of the product $\Gamma^0\Gamma^1\cdots\Gamma^{n-1}$, which is $\pm$ a fixed phase; equivalently, the last gamma matrix is $\pm$ a fixed multiple of the product of the first $n - 1$. \emph{Provenance:} standard; \cite[App.~B]{Polchinski1998}, \cite[\S 3]{VanProeyen1999}.
\emph{Source:} J. Polchinski, String Theory, Vol. II: Superstring Theory and Beyond, Cambridge University Press (1998), Appendix B., App.~B; A. Van Proeyen, Tools for supersymmetry, Ann. Univ. Craiova Phys. 9 (1999) 1--48, arXiv:hep-th/9910030., \S 3.
\end{fact}

\begin{theorem}[Rule 2 in the Kaluza-Klein algebra]\label{c:rule2.rule-2-kaluza-klein-algebra}
In $A_5 = \Cl(1,4)$:
\begin{enumerate}
\item[(i)] the volume element $\omega = \gamma e_4$ is central with $\omega^2 = +1$; $A_5$ is type C, $A_5 = \varepsilon_+ A_5 \oplus \varepsilon_- A_5$;
\item[(ii)] each ideal is a copy of the spacetime algebra: $a \mapsto \varepsilon_\pm a$ is an injective homomorphism $\Cl(1,3) \to \varepsilon_\pm A_5$ onto the ideal, and under it the fifth generator becomes minus or plus the spacetime pseudoscalar,
\[
\varepsilon_+ e_4 = -\,\varepsilon_+\gamma, \qquad \varepsilon_- e_4 = +\,\varepsilon_-\gamma :
\]
the two ideals are the two inequivalent five-dimensional Dirac representations of Fact \ref{c:rule2.imported-gamma-matrices-odd-dimension}, $\Gamma^4 = \mp\,\Gamma^0\Gamma^1\Gamma^2\Gamma^3$, on the one Dirac spinor of $\Cl(1,3) \cong \Mat_2(\HH)$;
\item[(iii)] the reflection of the fifth direction, $R: e_4 \mapsto -e_4$ with $e_0, \ldots, e_3$ fixed, extends to the automorphism $\hat{\cdot} \circ \Ad_{e_4}$ of $A_5$, sends $\omega \mapsto -\omega$, exchanges the ideals, and is outer: the five-dimensional parity is the exchange that Fact \ref{c:rule2.imported-classification-centre-type-c} says no inner automorphism can perform;
\item[(iv)] the Rule 2 repairs are $\Cl(2,4)$ ($\nu^2 = +1$) and $\Cl(1,5)$ ($\nu^2 = -1$), both $\cong \Mat_2(\Cl^0(1,4)) \cong \Mat_2(\Cl(1,3)) \cong \Mat_4(\HH)$, simple with centre $\R$, in which $\varepsilon_\pm$ are the diagonal matrix units;
\item[(v)] in either repair $\Ad_\nu$ restricts to the grade involution of $A_5$ and $\Ad_{\nu e_4}$ restricts to the parity $R$: the outer exchange has become inner, and it is a boost, $(\nu e_4)^2 = +1$, in $\Cl(2,4)$ and a rotation, $(\nu e_4)^2 = -1$, in $\Cl(1,5)$;
\item[(vi)] $\Cl(2,4)$ is the Rule 1 ascent of $\Cl(1,3)$, the conformal algebra of Minkowski space whose spinors carry the conformal spinor space $\CS$ (twistor space) (Fact \ref{c:rule2.imported-completion-spacetime-algebra-spinors}): the pair $(\nu, e_4)$ is its hyperbolic plane and $f = \tfrac12(\nu + e_4)$, $f^* = \tfrac12(\nu - e_4)$ its ladder pair; $\Cl(1,5)$ is the same algebra at residue $4$ with a definite adjoined plane. The residues run $6 \to 5 \to 6$ or $6 \to 5 \to 4$.
\end{enumerate}
\end{theorem}
\begin{proof}
(i) $\omega^2 = (-1)^{10}(-1)^4 = +1$ by Fact \ref{c:rule2.imported-classification-centre-type-c}, and $p - q = -3 \equiv 1 \pmod 4$; $\gamma e_4 = e_0e_1e_2e_3e_4 = \omega$ by definition. (ii) The map $a \mapsto \varepsilon_+ a$ is a homomorphism because $\varepsilon_+$ is a central idempotent. Injectivity: $\varepsilon_+ a = 0$ means $a = \omega a$; but for $a \in \Cl(1,3)$ every blade of $a$ omits $e_4$ while every blade of $\omega a = \gamma e_4 a$ contains it, so $a = \omega a$ forces $a = 0$. Both $\Cl(1,3)$ and the ideal have dimension $16$ (the ideal is $\Cl^0(1,4)$ by Fact \ref{c:rule2.imported-rule-2-matrix-units-forgotten}), so the map is onto. For the fifth generator, $e_4 = \gamma^{-1}\omega = -\gamma\omega$, so $\varepsilon_+ e_4 = -\gamma\varepsilon_+\omega = -\gamma\varepsilon_+$ since $\omega\varepsilon_+ = \varepsilon_+$, and $\varepsilon_- e_4 = -\gamma\varepsilon_-\omega = +\gamma\varepsilon_-$ since $\omega\varepsilon_- = -\varepsilon_-$. Fact \ref{c:rule2.imported-gamma-matrices-odd-dimension} identifies the two assignments $\Gamma^4 = \mp\Gamma^0\Gamma^1\Gamma^2\Gamma^3$ as the two inequivalent representations, and $\Cl(1,3) \cong \Mat_2(\HH)$ (\ref{c:std-algebra.sta-matrix}) has the one irreducible module $\HH^2$, the Dirac spinor. (iii) $\Ad_{e_4}$ negates $e_0, \ldots, e_3$ (each anticommutes with $e_4$, \ref{c:std-algebra.orthogonal-anticommute}) and fixes $e_4$; the grade involution negates all five; the composite fixes $e_0, \ldots, e_3$ and negates $e_4$, which is $R$ on generators, hence on the algebra. It negates $\omega$ (one generator negated) and so exchanges $\varepsilon_\pm$; an inner automorphism fixes the centre pointwise, so $R$ is outer. (iv) is Fact \ref{c:rule2.imported-rule-2-matrix-units-forgotten} with Lemma \ref{c:rule2.even-subalgebra-one-dimensional-extension} for $\Cl^0(1,4) \cong \Cl(1,3)$, and $\Mat_2(\Mat_2(\HH)) = \Mat_4(\HH)$. (v) $\Ad_\nu = \hat{\cdot}$ on $A_5$ is Fact \ref{c:rule2.imported-rule-2-matrix-units-forgotten}, and $\Ad_{\nu e_4} = \Ad_\nu \circ \Ad_{e_4} = \hat{\cdot} \circ \Ad_{e_4} = R$ on $A_5$ by (iii). $(\nu e_4)^2 = -\nu^2 e_4^2 = \nu^2$, which is $+1$ for $\Cl(2,4)$ and $-1$ for $\Cl(1,5)$; a bivector of square $+1$ generates a boost, of square $-1$ a rotation. (vi) is Theorem \ref{c:rule2.where-obstruction-sits-commuting-square}(c), (d) at $(p,q) = (1,3)$, with Fact \ref{c:rule2.imported-completion-spacetime-algebra-spinors} naming the algebra. The script checks every item: centrality and square of $\omega$, the idempotent identities, the rank $16$ of $a \mapsto \varepsilon_\pm a$ on the sixteen blades of $\Cl(1,3)$ and of each ideal, the two identities for $\varepsilon_\pm e_4$, $R = \hat{\cdot} \circ \Ad_{e_4}$ on generators, the matrix-unit relations and the trivial centre in both repairs, $\Ad_\nu$ and $\Ad_{\nu e_4}$ on generators in both, the two squares of $\nu e_4$, the ladder relations in $\Cl(2,4)$, the definiteness in $\Cl(1,5)$, the residues, and the Rule 1 corner $P\,\Cl(2,4)\,P$ of dimension $16$ commuting with $\Cl(1,3)$.

\end{proof}

\begin{fact}[Imported: the five-dimensional Chern-Simons shift]\label{c:rule2.imported-five-dimensional-chern-simons-shift}
A five-dimensional Dirac fermion of mass $m$ charged under a gauge field, integrated out, shifts the coefficient of the five-dimensional Chern-Simons term by $\tfrac12\sgn(m)$ (times the representation's cubic index for a non-abelian field); the sign is odd under the reflection of one spatial direction, so that the parity-violating term is present for a single species and absent for a pair whose contributions cancel. \emph{Provenance:} \cite{Witten1996, IMS1997}.
\emph{Source:} E. Witten, Phase transitions in M-theory and F-theory, Nucl. Phys. B 471 (1996) 195--216.; K. Intriligator, D. R. Morrison, N. Seiberg, Five-dimensional supersymmetric gauge theories and degenerations of Calabi-Yau spaces, Nucl. Phys. B 497 (1997) 56--100..
\end{fact}

\begin{remark}[What the Kaluza-Klein obstruction is, physically]\label{c:note.rule2.what-kaluza-klein-obstruction-physically}
Theorem \ref{c:rule2.rule-2-kaluza-klein-algebra}(ii) says that a five-dimensional fermion cannot avoid choosing an ideal: to write its Dirac operator one picks $\Gamma^4 = -\gamma$ or $+\gamma$, and no change of basis undoes the choice. Theorem \ref{c:rule2.rule-2-kaluza-klein-algebra}(iii) says that the five-dimensional parity exchanges the choices; Fact \ref{c:rule2.imported-five-dimensional-chern-simons-shift} says that the choice, together with the sign of the mass, is physical, carried by the induced Chern-Simons level. The obstruction is therefore not a formality of the real classification table: it is the statement that a five-dimensional fermion with a mass is parity-violating in a way its four-dimensional restriction is not, and that the violation is labelled by the ideal. Theorem \ref{c:rule2.rule-2-kaluza-klein-algebra}(iv)-(v) says what repairs it: a sixth direction, in which the two species become the $\varepsilon_+$ and $\varepsilon_-$ halves of one $\Mat_4(\HH)$-module, exchanged by the inner $\Ad_{\nu e_4}$. That is the algebraic form of the familiar statement that a six-dimensional spinor reduces to a five-dimensional fermion in each of the two representations, and that the pair is parity-invariant. Which sixth direction is Remark \ref{c:note.rule2.forgotten-bit-read}.
\end{remark}

\begin{remark}[The forgotten bit, read]\label{c:note.rule2.forgotten-bit-read}
Fact \ref{c:rule2.imported-rule-2-matrix-units-forgotten} says the fused algebra forgets the sign of $\nu^2$. Here the bit has a name on each side. With $\nu^2 = +1$ the repair is $\Cl(2,4)$: the sixth direction is a second time, the pair $(\nu, e_4)$ is the hyperbolic plane of the conformal completion, the parity of the fifth direction is a boost in that plane, and the fused algebra is the one whose spinors are the conformal spinors $\CS$ (twistors) (Fact \ref{c:rule2.imported-completion-spacetime-algebra-spinors}); the Kaluza-Klein direction and the conformal completion's null pair are the same two dimensions, differently sliced. With $\nu^2 = -1$ the repair is $\Cl(1,5)$: the sixth direction is a sixth space, the parity is a rotation into it, and the algebra sits at residue $4$. The two are isomorphic as algebras and the algebra cannot tell them apart; the geometry can, by whether the adjoined plane carries a null vector. The series' physical lineage has always chosen $\Cl(2,4)$, by the vector-level minimality of the root paper; the note observes that the same choice is, from the Kaluza-Klein side, the choice that returns the residue to $6$ and reads the fifth direction as half of a null pair.
\end{remark}

\begin{remark}[The conformal spinor space as a five-dimensional spinor]\label{c:note.rule2.twistor-five-dimensional-spinor}
The Fock space of Fact \ref{c:rule2.imported-completion-spacetime-algebra-spinors} is the complex spinor of $\Cl(2,4)$; as a module over $\Cl(1,4)_\C = \Mat_4(\C) \oplus \Mat_4(\C)$ it is $\varepsilon_+ \oplus \varepsilon_-$ halves, one copy of each inequivalent five-dimensional representation, since the odd element $\nu$ maps each half onto the other (Fact \ref{c:rule2.imported-rule-2-matrix-units-forgotten}); the two halves restrict to the same module over $\Cl^0(1,4) \cong \Cl(1,3)$, the Dirac spinor. The chirality halves of the six-dimensional spinor are a different splitting, by the six-dimensional volume element, which the odd elements of $\Cl(1,4)$ do not preserve; each chirality half restricted to $\mathrm{Spin}(1,4) \subset \mathrm{Spin}(2,4)$ is the five-dimensional Dirac spinor, the fundamental $\HH^2 = \C^4$ of $Sp(1,1) \subset SU(2,2)$. So the conformal spinor space $\CS$ (twistor space), as a representation of the five-dimensional Lorentz group (which is also the four-dimensional de Sitter group), is a five-dimensional Dirac spinor. This is standard group theory and the script does not check it; it is recorded because it is the spinor-side form of the square of a figure of the original.
\end{remark}

\section{Planar fermions}

\phantomsection\label{c:note.rule2.setting-planar-fermions}
The second system is the oldest home of the parity anomaly. Take generators $e_x, e_y$ of square $+1$ and $e_t$ of square $-1$: the algebra $\Cl(2,1)$ of a $2{+}1$-dimensional spacetime with time of square $-1$, which is the mirror's labelling in the series' convention (Corollary \ref{c:rule2.odd-lorentzian-dimension-exactly-one-real} at $n = 3$: the time-positive form $\Cl(1,2)$ is complex type).

\begin{fact}[Imported: the parity anomaly]\label{c:rule2.imported-parity-anomaly}
A single two-component Dirac fermion of mass $m$ in $2{+}1$ dimensions, coupled to a $U(1)$ gauge field, has an effective action containing a Chern-Simons term at level $\tfrac12\sgn(m)$ up to a sign that flips between the two inequivalent two-component representations; the term is odd under the reflection of one spatial axis, so a single species breaks parity, while two species, one in each representation with equal mass (equivalently the four-component fermion with its parity-invariant mass term), induce no net Chern-Simons term. \emph{Provenance:} \cite{NiemiSemenoff1983, Redlich1984, AGW1984}.
\emph{Source:} A. J. Niemi, G. W. Semenoff, Axial-anomaly-induced fermion fractionization and effective gauge-theory actions in odd-dimensional space-times, Phys. Rev. Lett. 51 (1983) 2077--2080.; A. N. Redlich, Gauge noninvariance and parity nonconservation of three-dimensional fermions, Phys. Rev. Lett. 52 (1984) 18--21; Parity violation and gauge noninvariance of the effective gauge field action in three dimensions, Phys. Rev. D 29 (1984) 2366--2374.; L. Alvarez-Gaum'e, E. Witten, Gravitational anomalies, Nucl. Phys. B 234 (1984) 269--330..
\end{fact}

\begin{fact}[Imported: graphene's two valleys, and the doubling]\label{c:rule2.imported-graphenes-two-valleys-doubling}
The low-energy theory of electrons on the honeycomb lattice consists of two two-component Dirac fermions, one at each of the two inequivalent corners $K, K'$ of the Brillouin zone, whose gamma matrices differ by the sign of one spatial gamma matrix; a lattice regularisation of a single two-component species is excluded by the fermion-doubling theorem, so the pair is forced. \emph{Provenance:} \cite{Semenoff1984, NielsenNinomiya1981, GSC2007}.
\emph{Source:} G. W. Semenoff, Condensed-matter simulation of a three-dimensional anomaly, Phys. Rev. Lett. 53 (1984) 2449--2452.; H. B. Nielsen, M. Ninomiya, Absence of neutrinos on a lattice (I). Proof by homotopy theory, Nucl. Phys. B 185 (1981) 20--40.; V. P. Gusynin, S. G. Sharapov, J. P. Carbotte, AC conductivity of graphene: from tight-binding model to 2+1-dimensional quantum electrodynamics, Int. J. Mod. Phys. B 21 (2007) 4611--4658..
\end{fact}

\begin{proposition}[Rule 2 in the plane]\label{c:rule2.rule-2-plane}
In $A_3 = \Cl(2,1)$:
\begin{enumerate}
\item[(i)] $\omega = e_xe_ye_t$ is central with $\omega^2 = +1$; $A_3 = \varepsilon_+ A_3 \oplus \varepsilon_- A_3$ with each ideal a copy of $\Cl^0(2,1) \cong \Cl(2,0) \cong \Mat_2(\R)$, and $\varepsilon_+ e_t = -\varepsilon_+ e_xe_y$, $\varepsilon_- e_t = +\varepsilon_- e_xe_y$: the two ideals are the two inequivalent two-component representations, which fix the time gamma matrix as $\mp$ the spatial bivector; the reflection of one spatial axis is $\hat{\cdot}\circ\Ad_{e_x}$, negates $\omega$, exchanges the ideals, and is outer;
\item[(ii)] the Rule 2 repairs are $\Cl(3,1)$ ($\nu^2 = +1$, a third spatial direction) and $\Cl(2,2)$ ($\nu^2 = -1$, a second time), both $\cong \Mat_2(\Mat_2(\R)) = \Mat_4(\R)$, in which $\Ad_\nu$ is the grade involution of $A_3$ and $\Ad_{\nu e_x}$ is the planar reflection: a rotation into the third direction, $(\nu e_x)^2 = -1$, in $\Cl(3,1)$, and a boost, $(\nu e_x)^2 = +1$, in $\Cl(2,2)$; $\Cl(3,1)$ is the Rule 1 ascent of $\Cl(2,0)$ and the residues run $2 \to 1 \to 2$ or $2 \to 1 \to 0$;
\item[(iii)] the time-positive form $\Cl(1,2)$ has central volume element of square $-1$: complex type, $\Cl(1,2) \cong \Mat_2(\C)$, no real splitting; the splitting appears on complexification.
\end{enumerate}
\end{proposition}
\begin{proof}
(i) $\omega^2 = (-1)^3(-1)^1 = +1$ and $p - q = 1$: type C by Fact \ref{c:rule2.imported-classification-centre-type-c}; the corner is $\Cl^0(2,1)$ by Fact \ref{c:rule2.imported-rule-2-matrix-units-forgotten} and $\Cl(2,0)$ by Lemma \ref{c:rule2.even-subalgebra-one-dimensional-extension}, which is $\Mat_2(\R)$ by \ref{c:std-algebra.cl11-matrix}. $e_t = (e_xe_y)^{-1}\omega = -e_xe_y\omega$ since $(e_xe_y)^2 = -1$, so $\varepsilon_\pm e_t = -e_xe_y\varepsilon_\pm\omega = \mp e_xe_y\varepsilon_\pm$. The reflection along $e_x$ is $\hat{\cdot}\circ\Ad_{e_x}$ by the argument of Theorem \ref{c:rule2.rule-2-kaluza-klein-algebra}(iii) with $e_x$ in place of $e_4$, and is outer for the same reason. (ii) Fact \ref{c:rule2.imported-rule-2-matrix-units-forgotten}, Lemma \ref{c:rule2.even-subalgebra-one-dimensional-extension}, and the computations of Theorem \ref{c:rule2.rule-2-kaluza-klein-algebra}(v)-(vi) with $\nu^2 = +1$ now the spacelike case: $(\nu e_x)^2 = -\nu^2e_x^2 = -\nu^2$. Theorem \ref{c:rule2.where-obstruction-sits-commuting-square}(c) at $(p,q) = (2,0)$, residue $2 \equiv 2 \pmod 4$, whose obstructing extension is the one by a generator of square $-1$, here $e_t$, and whose opposite-sign repair is $\nu^2 = +1$, gives $\Cl(3,1)$ as the Rule 1 ascent of $\Cl(2,0)$. (iii) $\omega^2 = (-1)^3(-1)^2 = -1$ for $\Cl(1,2)$, $p - q = -1 \equiv 3 \pmod 4$; the isomorphism with $\Mat_2(\C)$ is the classification, \ref{c:primer.classification-real-clifford-algebras}. The script checks (i)-(iii) in full, the residues included.

\end{proof}

\begin{remark}[What the planar obstruction is, physically]\label{c:note.rule2.what-planar-obstruction-physically}
Fact \ref{c:rule2.imported-parity-anomaly} is Proposition \ref{c:rule2.rule-2-plane}(i) with a coupling switched on: the ideal a two-component fermion lives in is the sign of its induced Chern-Simons level, the parity that exchanges the ideals is the reflection that flips the level, and a single species cannot be parity-invariant because no inner change of variables moves it to the other ideal. The four-component fermion of planar electrodynamics is Proposition \ref{c:rule2.rule-2-plane}(ii): a module over $\Mat_4(\R)$, the fused algebra, on which the two species are the $\varepsilon_\pm$ halves and the matrix $\nu = \gamma^3$, which anticommutes with the three planar gamma matrices, is the inner automorphism that exchanges them; the parity-invariant mass is the one that respects the fusion. Fact \ref{c:rule2.imported-graphenes-two-valleys-doubling} is the same proposition realised in a crystal: the two valleys are the two ideals, related by the sign of one spatial gamma matrix, which by (i) is exactly the reflection that exchanges them, and the doubling theorem is the statement that a lattice, having a third direction available to it in the form of a lattice regulator, cannot host the unfused algebra. In the series' language: nature has done the fusion, and the two-component fermion is what one sees when one looks at a single ideal of it.
\end{remark}

\begin{remark}[Eleven dimensions, and the reality of spinors]\label{c:note.rule2.eleven-dimensions-reality-spinors}
Corollary \ref{c:rule2.odd-lorentzian-dimension-exactly-one-real} at $n = 11$ says $\Cl(10,1) \cong \Mat_{32}(\R) \oplus \Mat_{32}(\R)$ is split while $\Cl(1,10) \cong \Mat_{32}(\C)$ has complex type. The irreducible real module of the first is $\R^{32}$ and of the second is $\C^{32}$, which is the classical statement that Majorana spinors exist in eleven dimensions with time of square $-1$ and not with time of square $+1$, the reason eleven-dimensional supergravity is written in the former convention. \emph{Provenance:} \cite{KugoTownsend1983, VanProeyen1999}. The note records this as a third physical face of the same theorem, without a section of its own: the Rule 2 repair of $\Cl(10,1)$ is $\Cl(11,1)$ or $\Cl(10,2)$, both $\Mat_{64}(\R)$; the former is the Rule 1 ascent of $\Cl(10,0)$, and the latter is the algebra of the two-time proposals, which the note mentions and does not endorse.
\end{remark}

\section{What has been exercised, and what has not}

\phantomsection\label{c:note.rule2.setting-what-has-been-exercised-what}
The ledger asked for a physical system whose natural algebra carries a central splitting. The note has given two, and the theorem that they are instances of: the five-dimensional Kaluza-Klein algebra of this world and the planar algebra of the mirror's convention, each split, each with its ideals the two inequivalent spinor representations that physics has always known about, each with its outer exchange the parity that physics has always known to be anomalous, each with its Rule 2 repair the doubled fermion that physics has always used to restore the symmetry. In the five-dimensional case the repair is the conformal algebra of the series' own root paper, and Theorem \ref{c:rule2.where-obstruction-sits-commuting-square} says that this is not a coincidence of the example but the general shape: for the physical lineage, the hyperbolic plane that Rule 1 adjoins is a Kaluza-Klein direction plus the direction that repairs it. Rule 2 has been exercised physically in the sense that the obstruction, the exchange, and the repair are each carried by a measured or standard physical structure, imported here as Facts.

What has not been exercised is dynamics. The note does not derive that a sixth direction exists, nor that the fifth one does; it shows that if a fifth is adjoined the algebra obstructs, and that the repair is the direction the series already had. The forgotten bit is read, not decided: whether the parity of the fifth direction becomes a boost or a rotation is a geometric datum, and the series' choice of $\Cl(2,4)$ rests on the root paper's minimality argument, not on anything here. The planar system exercises the fusion at the level of representation theory and of an induced term in an effective action; it does not make the third direction physical, which in graphene it is not. And the note has used the series' convention where it needed to: the Kaluza-Klein algebra splits in the time-positive form, the planar and eleven-dimensional algebras in the time-negative one, and the complexified statement of Corollary \ref{c:rule2.odd-lorentzian-dimension-exactly-one-real} is the convention-free residue.

The row's new state, for the capstone's ledger: owned by this note; the kinematic exercise delivered; remaining, a system in which the fused direction is forced dynamically rather than by the representation theory of the fermions, which is the question of what physical principle selects the sign of $\nu^2$.

\section{Position in the series}

\phantomsection\label{c:note.rule2.setting-position-series}
This note is a companion to T17 of the series, the capstone \cite{ButlerT17}, and closes Direction 7 of its forward ledger. Every algebraic input is the root paper T1 \cite{ButlerT1}: the classification and the centre (its Theorem 2.5 and Lemmas 2.6, 2.7), the existence and outerness of type C (Theorem 3.6), the Rule 2 machinery (Proposition 5.2, Lemmas 5.3 and 5.4, Theorems 5.5 and 5.7, Remark 5.8), Rule 1's Morita triviality (Theorem 7.1), and the parking of the clock (Proposition 7.4); the note's Lemma \ref{c:rule2.even-subalgebra-one-dimensional-extension} is the one standard fact it proves for itself. Theorem \ref{c:rule2.where-obstruction-sits-commuting-square} is the note's contribution to T1's own subject: the two rules, compared there as amplifications of the whole algebra and of an even half, coincide on every one-dimensional extension of an even-dimensional algebra, where the even half is the base. The naming of $\Cl(2,4)$ as the conformal completion and of its spinors as twistors is the root paper's and the twistor paper T9's \cite{ButlerT9}. The residue set $\{2, 6\}$ is the dark-matter paper T14's \cite{ButlerT14}, whose reading of the signature convention as indexical is what Corollary \ref{c:rule2.odd-lorentzian-dimension-exactly-one-real} needs. The primer T0 \cite{ButlerT0} carries the classification table in the form the note quotes. Nothing downstream depends on this note; the capstone's ledger row is its only consumer.

\section{Machine verification}

\phantomsection\label{c:note.rule2.setting-machine-verification}
Every computational identity asserted in this note has been verified in exact rational arithmetic by the script \texttt{verify\_T17\_rule2.py}, which imports the series' shared library \texttt{GA.py}: its basis-blade Clifford engine, exact ranks, and exact congruence signatures; no floating point anywhere. Twenty-four checks, all passing, in six parts.

\emph{Part A: the extension theorem.} On every even-dimensional nondegenerate $(p,q)$ with $2 \leq p + q \leq 6$, the volume element of each one-dimensional extension is tested for centrality and its square computed: the spacelike extension is type C exactly when $p - q \equiv 2 \pmod 4$, the timelike one exactly when $p - q \equiv 0$, the other has complex type, and the residue sets are $\{2, 6\}$ and $\{0, 4\}$. The Lorentzian alternation with its complex-type complement in dimensions $3, 5, 7, 9, 11$, and the three named cases.

\emph{Part B: the even subalgebra.} The map $e_i \mapsto e_i\nu$ on $(1,3)$ and $(2,0)$ with both signs of $\nu^2$: relations, evenness, and the generated dimension $2^n$.

\emph{Part C: the Kaluza-Klein algebra.} Centrality and square of $\omega$, the idempotent identities, $\gamma^2 = -1$ and $\omega = \gamma e_4$; rank $16$ of $a \mapsto \varepsilon_\pm a$ on the blades of $\Cl(1,3)$ and of each ideal; $\varepsilon_\pm e_4 = \mp\varepsilon_\pm\gamma$; the reflection as $\hat{\cdot}\circ\Ad_{e_4}$, its action on $\omega$ and on $\varepsilon_\pm$.

\emph{Part D: the two fusions.} In $\Cl(2,4)$ and $\Cl(1,5)$: the sixteen matrix-unit products, the non-centrality of $\varepsilon_\pm$, the trivial centre; $\Ad_\nu$ and $\Ad_{\nu e_4}$ on every generator; $(\nu e_4)^2 = \pm 1$; the ladder relations of $(\nu, e_4)$ in $\Cl(2,4)$ and the signature $(0,2)$ in $\Cl(1,5)$; the residues; the Rule 1 projectors $ff^*$, $f^*f$ commuting with $\Cl(1,3)$ and the corner of dimension $16$.

\emph{Part E: the plane.} The same programme for $\Cl(2,1)$: the splitting, the ideal dimension $4$, $\varepsilon_\pm e_t = \mp\varepsilon_\pm e_xe_y$, both fusions with their matrix units, inner grade involution and inner reflection, $(\nu e_x)^2 = \mp 1$, the complex type of $\Cl(1,2)$ and of $\Cl(4,1)$, and the residues.

\emph{Part F: the square in general.} On the grid of Part A: the obstructed extension is type C, the opposite-sign repair is $\Cl(p{+}1,q{+}1)$ at the same residue with a hyperbolic adjoined plane carrying a ladder pair, the same-sign repair sits two residues away with a definite plane, and every repaired algebra has trivial centre.

What the script does not check: the physics imported as Facts \ref{c:rule2.imported-gamma-matrices-odd-dimension}, \ref{c:rule2.imported-five-dimensional-chern-simons-shift}, \ref{c:rule2.imported-parity-anomaly}, and \ref{c:rule2.imported-graphenes-two-valleys-doubling} (the note identifies structures, it does not recompute effective actions); the isomorphism $\Cl(1,3) \cong \Mat_2(\HH)$ and the rest of the classification table beyond the dimension counts and centres it computes (imported); the group-theoretic restriction of Remark \ref{c:note.rule2.twistor-five-dimensional-spinor} (standard); and the eleven-dimensional module counts of Remark \ref{c:note.rule2.eleven-dimensions-reality-spinors} beyond the split-versus-complex verdict of Part A.

\section{Results used from elsewhere}
The proofs above use the following results of the series, stated briefly; each is proved in the paper named.

\begin{fact}[The Fock model of the complexified ascended algebra]\label{c:correspondence.fock-model-complexified-ascended-algebra}
The exterior algebra on a maximal isotropic subspace, with wedge and contraction operators, is the unique irreducible module of the complexified Clifford algebra $\Cl(2,4)_\C$, and its even and odd halves are the two four-dimensional chiralities.
\emph{Source:} T9, Thm 3.1 (\href{https://doi.org/10.5281/zenodo.22722290}{doi:10.5281/zenodo.22722290}).
\end{fact}

\begin{fact}[The spinor identification]\label{c:correspondence.spinor-identification}
There is a canonical $\Spin$-equivariant isomorphism from $V_\C$ to the bivectors of the conformal spinor space carrying the quadratic form to the wedge square, so null vectors are exactly the simple bivectors, that is, the lines of the projective spinor space.
\emph{Source:} T9, Thm 3.4 (\href{https://doi.org/10.5281/zenodo.22722290}{doi:10.5281/zenodo.22722290}).
\end{fact}

\begin{fact}[Classification of the real Clifford algebras]\label{c:primer.classification-real-clifford-algebras}
$\Cl(p,q)$ is a matrix algebra over $\R$, $\C$ or $\HH$, or a sum of two such, according to $p - q \bmod 8$, and it is simple except when $p - q \equiv 1 \pmod 4$.
\emph{Source:} T0, Thm 9.3 (\href{https://doi.org/10.5281/zenodo.22714501}{doi:10.5281/zenodo.22714501}).
\end{fact}

\begin{fact}[Cl(1,1) and Cl(2,0) are the real 2x2 matrices]\label{c:std-algebra.cl11-matrix}
$\Cl(H) \cong \Mat_2(\R)$ for a hyperbolic plane, with a hyperbolic pair going to the off-diagonal matrix units; also $\Cl(2,0) \cong \Mat_2(\R)$.
\emph{Source:} the algebra core of the Transport Lab, proved there with an exact check; standard, see P. Lounesto, Clifford Algebras and Spinors, 2nd ed., London Mathematical Society Lecture Note Series 286, Cambridge University Press (2001), Ch. 16 (matrix representations and periodicity of 8).
\end{fact}

\begin{fact}[Complexification, and the odd-dimensional split]\label{c:std-algebra.complexification}
$\Cl(V) \otimes \C \cong \Cl_n(\C)$, a complex matrix algebra (doubled for odd $n$); for odd $n$, $\Cl \cong \Cl^{+} \oplus \Cl^{+}$ if $\omega^{2} = 1$ and $\Cl \cong \Cl^{+} \otimes \C$ if $\omega^{2} = -1$.
\emph{Source:} the algebra core of the Transport Lab, proved there with an exact check; standard, see P. Lounesto, Clifford Algebras and Spinors, 2nd ed., London Mathematical Society Lecture Note Series 286, Cambridge University Press (2001), Ch. 16 (matrix representations and periodicity of 8).
\end{fact}

\begin{fact}[The Clifford algebra exists and is unique]\label{c:std-algebra.existence-uniqueness}
$\Cl(V, Q) = T(V)/\langle v \otimes v - Q(v) \rangle$ exists and is unique up to unique isomorphism; the wedge-plus-contraction operators on $\Lambda V$ give a concrete representation.
\emph{Source:} the algebra core of the Transport Lab, proved there with an exact check; standard, see C. Chevalley, The Algebraic Theory of Spinors, Columbia University Press (1954), Ch. II (the Clifford algebra of a quadratic space).
\end{fact}

\begin{fact}[Orthogonal vectors anticommute]\label{c:std-algebra.orthogonal-anticommute}
$uv = -vu$ iff $B(u,v) = 0$; in an orthogonal basis $e_I e_J = \pm \prod_{I \cap J} Q(e_i)\; e_{I \triangle J}$ with the sorting sign.
\emph{Source:} the algebra core of the Transport Lab, proved there with an exact check; standard, see C. Doran and A. Lasenby, Geometric Algebra for Physicists, Cambridge University Press (2003), Ch. 4 (foundations of geometric algebra).
\end{fact}

\begin{fact}[The spacetime algebra as a matrix algebra, and the spacetime split]\label{c:std-algebra.sta-matrix}
$\Cl(1,3) \cong \Mat_2(\HH)$, the Dirac matrices represent $\C \otimes \Cl(1,3) \cong \Mat_4(\C)$, and $\sigma_k = \gamma_k\gamma_0$ give $\Cl^{+}(1,3) \cong \Cl(3,0)$ with $\sigma_1\sigma_2\sigma_3 = I$.
\emph{Source:} the algebra core of the Transport Lab, proved there with an exact check; standard, see C. Doran and A. Lasenby, Geometric Algebra for Physicists, Cambridge University Press (2003), Ch. 5 (relativity and spacetime).
\end{fact}

\begin{fact}[The ascent renders the grade involution inner]\label{c:transport.ascent-renders-grade-involution-inner}
In the one-generator ascent, conjugation by the adjoined generator $\nu$ is the grade involution on $A$, exchanging the projectors $\varepsilon_\pm$, so either one generates the whole algebra.
\emph{Source:} T1, Prop 5.2 (\href{https://doi.org/10.5281/zenodo.21743863}{doi:10.5281/zenodo.21743863}).
\end{fact}

\begin{fact}[The centre]\label{c:transport.centre}
For $n = p + q \geq 1$ the centre of $\Cl(p,q)$ is $\R \cdot 1$ when $n$ is even and $\R \cdot 1 \oplus \R \cdot \omega$ when $n$ is odd, with $\omega$ the pseudoscalar.
\emph{Source:} T1, Lem 2.6 (\href{https://doi.org/10.5281/zenodo.21743863}{doi:10.5281/zenodo.21743863}).
\end{fact}

\begin{fact}[Classification of the real Clifford algebras]\label{c:transport.classification-real-clifford-algebras}
The real Clifford algebra $\Cl(p,q)$ is a matrix algebra over $\R$, $\C$ or $\HH$, or a sum of two such, according to $p - q \bmod 8$, and it is simple except when $p - q \equiv 1 \pmod 4$.
\emph{Source:} M. F. Atiyah, R. Bott and A. Shapiro, Clifford modules (1964); I. R. Porteous, Clifford Algebras and the Classical Groups (1995); P. Lounesto, Clifford Algebras and Spinors (2001).
\end{fact}

\begin{fact}[The corner is the even subalgebra]\label{c:transport.corner-even-subalgebra}
For $p - q \equiv 1 \pmod 4$, $a \mapsto \varepsilon_+ a$ is a unital isomorphism $\Cl^{0}(p,q) \to \varepsilon_+ A$, and $\varepsilon_+ B_\gamma \varepsilon_+ = \varepsilon_+ A$.
\emph{Source:} T1, Lem 5.3 (\href{https://doi.org/10.5281/zenodo.21743863}{doi:10.5281/zenodo.21743863}).
\end{fact}

\begin{fact}[Existence and classification of type C]\label{c:transport.existence-classification-type-c}
$\Cl(p,q)$, $n \geq 1$, has a nontrivial central idempotent exactly when $p - q \equiv 1 \pmod 4$, and then the only ones are $\varepsilon_\pm = \tfrac12(1 \pm \omega)$, which split the algebra into two simple ideals and are exchanged by the grade involution but by no inner automorphism.
\emph{Source:} T1, Thm 3.6 (\href{https://doi.org/10.5281/zenodo.21743863}{doi:10.5281/zenodo.21743863}).
\end{fact}

\begin{fact}[Fusion occurs at most once]\label{c:transport.fusion-occurs-most-once}
Along any transport history Rule 1 fixes the mod-8 residue and Rule 2 leaves the type-C residues for good, so fusion happens at most once and the Wedderburn type is then fixed.
\emph{Source:} T1, Prop 7.4 (\href{https://doi.org/10.5281/zenodo.21743863}{doi:10.5281/zenodo.21743863}).
\end{fact}

\begin{fact}[Matrix units, and the Peirce isomorphism]\label{c:transport.matrix-units-peirce-isomorphism}
The projectors $\varepsilon_\pm$ and the adjoined generator $\nu$ form a full system of $2 \times 2$ matrix units, so $B_\gamma \cong \Mat_2(E_{11} B_\gamma E_{11})$.
\emph{Source:} T1, Lem 5.4 (\href{https://doi.org/10.5281/zenodo.21743863}{doi:10.5281/zenodo.21743863}).
\end{fact}

\begin{fact}[Minimality]\label{c:transport.minimality}
Every null-lift completion adjoins directions of both signs, so the minimal completion is the hyperbolic-plane ascent $\Cl(p,q) \to \Cl(p{+}1,q{+}1)$, unique up to isometry, with increment $(+1,+1)$.
\emph{Source:} T1, Thm 4.13 (\href{https://doi.org/10.5281/zenodo.21743863}{doi:10.5281/zenodo.21743863}).
\end{fact}

\begin{fact}[No-go at the level of elements]\label{c:transport.no-go-level-elements}
In mixed signature with $n \geq 3$, no assignment of von Neumann pseudo-inverses to null vectors, valued in the Clifford algebra of any extension, is equivariant even under $\Spin^0(p,q)$.
\emph{Source:} T1, Thm 4.2 (\href{https://doi.org/10.5281/zenodo.21743863}{doi:10.5281/zenodo.21743863}).
\end{fact}

\begin{fact}[One generator is minimal; the minimal completions classified]\label{c:transport.one-generator-minimal-minimal-completions-classified}
For $p - q \equiv 1 \pmod 4$, the minimal simple Clifford completions of $\Cl(p,q)$ are exactly the two one-generator ascents $\Cl(p{+}1,q)$ and $\Cl(p,q{+}1)$, isomorphic as algebras, while a null generator only thickens the algebra.
\emph{Source:} T1, Thm 5.7 (\href{https://doi.org/10.5281/zenodo.21743863}{doi:10.5281/zenodo.21743863}).
\end{fact}

\begin{fact}[Rule 1 is Morita-trivial]\label{c:transport.rule-1-morita-trivial}
Rule 1 gives $\Cl(p+1, q+1) \cong \Mat_2(\Cl(p,q))$ with the adjoined ladder pair as off-diagonal matrix units, so it preserves the Morita class, the Wedderburn type and the clock residue $p - q \pmod 8$.
\emph{Source:} T1, Thm 7.1 (\href{https://doi.org/10.5281/zenodo.21743863}{doi:10.5281/zenodo.21743863}).
\end{fact}

\begin{fact}[Rule 2: ideal fusion, sharp form]\label{c:transport.rule-2-ideal-fusion-sharp-form}
For $p - q \equiv 1 \pmod 4$, adjoining one generator $\nu$ with $Q(\nu) = \pm 1$ gives the simple algebra $\Mat_2(\Cl^0(p,q))$ with centre $\R$, the same for either sign.
\emph{Source:} T1, Thm 5.5 (\href{https://doi.org/10.5281/zenodo.21743863}{doi:10.5281/zenodo.21743863}).
\end{fact}

\begin{fact}[The sign of $\omega^2$]\label{c:transport.sign}
The pseudoscalar of $\Cl(p,q)$ satisfies $\omega^2 = (-1)^{(p-q)(p-q-1)/2}$, so $\omega^2 = +1$ exactly when $p - q \equiv 0$ or $1 \pmod 4$.
\emph{Source:} T1, Lem 2.7 (\href{https://doi.org/10.5281/zenodo.21743863}{doi:10.5281/zenodo.21743863}).
\end{fact}

\begin{refutation}[The two one-generator ascents differ as algebras]\label{c:refuted.the-two-one-generator-ascents-differ-as-algebras}
Refuted: the two one-generator ascents are always isomorphic as algebras, the sign of the adjoined generator being a geometric bit the algebra forgets.
\end{refutation}

\begin{refutation}[The original conjecture, refuted below]\label{c:transport.original-conjecture-refuted-below}
Element-by-element repair of a null vector by an equivariant partner one ascent up: refuted.
\end{refutation}

\begin{thebibliography}{99}
\bibitem{ButlerT1} L.~Butler,
\emph{The Transport Principle: Total Division in Clifford Algebras via Minimal Completion},
Transport Papers T1 (2026).
\bibitem{ButlerT17} L.~Butler,
\emph{Capstone: The Transport Programme, Method, Status, and the Falsifiable Ledger},
Transport Papers T17 (2026).
\bibitem{ButlerT14} L.~Butler,
\emph{Dark Matter and the Arrow of Time: The Five Dark Towers},
Transport Papers T14 (2026).
\bibitem{ABS1964} M.~F.~Atiyah, R.~Bott and A.~Shapiro,
Clifford modules,
\emph{Topology} \textbf{3} (1964), suppl.~1, 3--38.
\bibitem{Lounesto2001} P.~Lounesto,
\emph{Clifford Algebras and Spinors}, 2nd ed.,
London Mathematical Society Lecture Note Series 286, Cambridge University Press (2001).
\bibitem{ButlerT9} L.~Butler,
\emph{The Twistor Correspondence: Ascended Spinors, Incidence, and the Palatial Bridge},
Transport Papers T9 (2026).
\bibitem{Polchinski1998} J.~Polchinski,
\emph{String Theory, Vol.~II: Superstring Theory and Beyond},
Cambridge University Press (1998), Appendix B.
\bibitem{VanProeyen1999} A.~Van Proeyen,
\emph{Tools for supersymmetry},
Ann.\ Univ.\ Craiova Phys.\ 9 (1999) 1--48, arXiv:hep-th/9910030.
\bibitem{Witten1996} E.~Witten,
\emph{Phase transitions in M-theory and F-theory},
Nucl.\ Phys.\ B 471 (1996) 195--216.
\bibitem{IMS1997} K.~Intriligator, D.~R. Morrison, N.~Seiberg,
\emph{Five-dimensional supersymmetric gauge theories and degenerations of Calabi-Yau spaces},
Nucl.\ Phys.\ B 497 (1997) 56--100.
\bibitem{NiemiSemenoff1983} A.~J. Niemi, G.~W. Semenoff,
\emph{Axial-anomaly-induced fermion fractionization and effective gauge-theory actions in odd-dimensional space-times},
Phys.\ Rev.\ Lett.\ 51 (1983) 2077--2080.
\bibitem{Redlich1984} A.~N. Redlich,
\emph{Gauge noninvariance and parity nonconservation of three-dimensional fermions},
Phys.\ Rev.\ Lett.\ 52 (1984) 18--21; \emph{Parity violation and gauge noninvariance of the effective gauge field action in three dimensions}, Phys.\ Rev.\ D 29 (1984) 2366--2374.
\bibitem{AGW1984} L.~Alvarez-Gaum\'e, E.~Witten,
\emph{Gravitational anomalies},
Nucl.\ Phys.\ B 234 (1984) 269--330.
\bibitem{Semenoff1984} G.~W. Semenoff,
\emph{Condensed-matter simulation of a three-dimensional anomaly},
Phys.\ Rev.\ Lett.\ 53 (1984) 2449--2452.
\bibitem{NielsenNinomiya1981} H.~B. Nielsen, M.~Ninomiya,
\emph{Absence of neutrinos on a lattice (I). Proof by homotopy theory},
Nucl.\ Phys.\ B 185 (1981) 20--40.
\bibitem{GSC2007} V.~P. Gusynin, S.~G. Sharapov, J.~P. Carbotte,
\emph{AC conductivity of graphene: from tight-binding model to 2+1-dimensional quantum electrodynamics},
Int.\ J.\ Mod.\ Phys.\ B 21 (2007) 4611--4658.
\bibitem{KugoTownsend1983} T.~Kugo, P.~Townsend,
\emph{Supersymmetry and the division algebras},
Nucl.\ Phys.\ B 221 (1983) 357--380.
\bibitem{ButlerT0} L.~Butler,
\emph{An Introduction to Geometric Algebra: A Primer for the Transport Papers},
Transport Papers T0 (2026).
\end{thebibliography}
\end{document}